\documentclass[12pt]{article}
\usepackage{graphicx} % Required for inserting images
\usepackage{float}
\usepackage{rotating, graphicx}
\usepackage{booktabs}
\usepackage{multirow}
\usepackage{longtable}
\usepackage{array}
\usepackage{placeins}
\usepackage{caption}
\usepackage{subcaption}
\linespread{1.25}
\usepackage[a4paper,left=2.54cm,right=2.54cm,top=2.54cm,bottom=2.54cm]{geometry}
\setlength{\parindent}{20pt
    \setlength{\parskip}{8pt plus 2pt minus 1pt}
}
\usepackage[natbibapa]{apacite}
\usepackage{multibib}
\usepackage{xcolor}
\usepackage{url}
\usepackage{etoolbox}
\usepackage{appendix}
\usepackage{amsmath}
\usepackage{mathrsfs}
\usepackage{amssymb}
\usepackage{etoc}
\usepackage{pdflscape}
\def\changemargin#1#2{\list{}{\rightmargin#2\leftmargin#1}\item[]}
\let\endchangemargin=\endlist %https://tex.stackexchange.com/a/600

\widowpenalties 4 10000 10000 10000 0
\clubpenalties 4 10000 10000 10000 0



\renewcommand{\bibfont}{\small} % Size of references
\renewcommand\bibsection{%
  \section*{\refname\markboth{\MakeUppercase{\refname}}{\MakeUppercase{\refname}}}%
  \addcontentsline{toc}{section}{\refname}%
}

\begin{document}

\begin{titlepage}

    \begin{center}
        \LARGE
        \textbf{A Perfect Storm: First-Nature Geography and Economic Development$^*$} \\
        
        \vspace{0.5cm}
        \large
        Christian Vedel, \\ University of Southern Denmark,\\
        \small
        \vspace{0.25cm}
        christian-vs@sam.sdu.dk; 
        \url{www.sites.google.com/view/christianvedel} 
        
        \vspace{0.75cm}
    
        \large
        \textbf{Abstract} \\     
    \end{center}

   
    \small
    \begin{changemargin}{1cm}{1cm}
    In 1825 a storm cut a new channel through Denmark’s Limfjord, providing an exogenous shock to first‑nature geography. Difference‑in‑differences estimates show the channel increased trade immediately and, within a generation, lifted population by 26.7 percent — an elasticity of 1.6 relative to the improved market access. Higher fertility and economic growth of new industries, not migration, drove the expansion. A mirror experiment—the waterway’s closure circa 1086‑1208 — caused symmetric declines in medieval coin and building finds, bolstering external validity. These results offer the first robust causal evidence that first‑nature geomorphology shapes the location of economic activity.
    
    \vspace{0.05cm} 
    \textbf{JEL codes}: N01, N73, O18, R1 \\
    \textbf{Keywords:} First-nature, Trade, Geography, Natural Experiment
    \end{changemargin}

    
    
    \vfill
    
    \footnotesize
    $^*$This project was made possible through generous funding by the Independent Research Fund Denmark (DFF – 6109-00123). I want to thank Paul Sharp, Christian Møller-Dahl, Casper Worm Hansen, Max Schulze, Gregory Clark, James Fenske, Mathias Barding, Torben Johansen, Nadja van't Hoff, Kerstin Enflo, Neil Cummins, and former PhD student-colleagues at SDU and at LSE for valuable feedback. This project also benefitted from research assistance of Andreas Slot Ravnholt. Comments from participants at workshops and conferences, including DGPE (2021, 2022), EHES (2022, Groningen), SSEH (2022, Gothenburg), LSE Economic History Graduate Seminar (2022), EHS (2023), University of Copenhagen End-of-Semester workshop (2023), Annual CITP (Nottingham), Danish Society of Transport Economics, and SDU PhD seminars are gratefully acknowledged. This paper also benefitted from improvements suggested by ChatGPT. Important early inspiration came from the podcast 'Kongerækken' and Poulsen (2019). Finally, I want to thank Maria Fay Courtney Bohr. On our vacation during the COVID-summer of 2020, we stumbled over some fascinating local history resulting in the present paper.  \\
    Code and replication materials: \url{https://github.com/christianvedels/A_perfect_storm}.
    

\end{titlepage}
\newpage

\section{Introduction} \label{intro}

% Geography–institutions puzzle; sets big question
Economic disparities between regions is a central focus of both academic and policy debates. While institutions, culture, and infrastructure investments are often highlighted as key determinants of prosperity, the role of geography remains an enduring and contentious issue. Are certain regions destined to prosper simply because of their geography? This is pertinent since geography is changing in the face of climate change.\footnote{\citet{kavan2025new, IPCCAR6SYR2023, Mentaschi2018}} Yet the extent to which first-nature geography — (otherwise) immutable geographic features such as rivers, coastlines, and mountains — shapes the location and level of economic development remain incompletely mapped in the empirical mist.

% Introduces 1825 storm natural experiment & causal 
While it is well established that geography matters \citep{Henderson2018satelite, Davis2002}, disentangling the causal influence of first-nature geography from institutions, culture, and subsequent infrastructure investments remains a formidable challenge. This paper addresses this by examining a rare instance of time‑varying first‑nature geography: a sudden new waterway created by a violent storm in 1825 in the West Limfjord region of north‑western Denmark (see Figure \ref{fig:main_map}). At the time, the West Limfjord region was one of the poorest and most isolated corners of Denmark. First-nature geography usually changes slowly over centuries. But here geography changed unexpectedly over night. The breach granted direct access to North Sea shipping lanes and the rapidly industrializing ports of Britain, slashing transport costs and widening market opportunities. This natural experiment allows an investigation into a central unanswered question: What is the causal effect of first‑nature geography on the location of economic activity?

% Method + headline results (DiD, trade ↑, pop ↑)
The rapid change in first-nature geography is analyzed using a difference-in-differences framework. Combined with newly available micro-level data for nineteenth-century Denmark, this approach allows estimation of the effect of first-nature geography on several measures of economic activity: trade volumes, population density, occupational structure, fertility, and internal migration. The newly formed channel led to an immediate increase in trade to and from the affected region, whereas population density and occupational structure adjusted more slowly, consistent with Malthusian dynamics. A key takeaway is that the consequences of changes in first-nature geography can take decades to materialize. Initial challenges gave way for long-run economic benefits, which I can observe in patterns of population growth stemming from higher fertility rather than internal migration. It was also accompanied by a larger share of residents working in manufacturing and fishing. Together, these changes point to higher long-run prosperity from favorable first-nature geography.

% Data sources & treatment definition detailed
This analysis draws on several data sources. For data on trade, the paper draws on records from the Sound-Toll-Registers \citep{soundtoll_data}, which contain detailed information on the origin and destination of every ship passing between the North Sea and the Baltic. From these registers, data for all Danish ports (126 of them) were extracted, and the count of ships to and from each port is used as a measure of trade. For other outcomes — such as population density, fertility, and occupational structure — micro‑level Danish census data \citep{mathiesen2022linklives} are used. This data set covers 1,589 parishes and standard demographic information as well as occupational descriptions, which are standardized into \textit{HISCO} \citep{leeuwen2002hisco} using a new language model for historical occupations \citep{dahl2024breaking}. All data are incorporated into a general difference‑in‑differences framework comparing \textit{affected} regions with the rest of Denmark. Affected regions are defined both by a simple dummy variable and by a continuous measure of the change in market access generated by the new waterway.

% Medieval channel closure gives external‑validity test
A central concern is whether these results generalize to other contexts. This is addressed directly. In a parallel — but reversed — natural experiment, a channel in the same location closed between 1086 and 1208. This setting provides strong external validity because medieval Denmark differed markedly in culture, religion, technology, and institutions. Because of this, this works well to test, whether effects are similar in a very different setting. Northwestern Denmark had enjoyed a westward channel, used by traders and Viking fleets since the last Ice Age; its closure, driven by gradual land rises \citep{Christensen2004}, led to long‑run stagnation in the West Limfjord until the 1825 reversal. Using archaeological data,\footnote{Which in itself contributes to a recent strand of literature using archaeological data to infer economic activity \citep{Barjamovic2019, Bakker2021Phonecians, Allen2023, boehm2024}.} it is shown that the sudden \textit{removal} of favorable first‑nature geography caused a 2 percent decline in coin findings.

% Positions contribution vs. first‑/second‑nature literature
The central contribution of this paper connects closely to \cite{Ahlfeldt2015}, who examine second‑nature geography through the lens of the Berlin Wall’s construction and demolition. The present study offers analogous causal evidence, but for first‑nature (rather than second‑nature) geography’s role in determining the location of prosperity. \cite{Henderson2018satelite} document that geography—via trade and agriculture — correlates strongly enough to explain 47 percent of the worldwide variation in modern economic activity, whereas \cite{Davis2002} argue that the location of economic activity results from a combination of “locational fundamentals” (first‑nature) and increasing returns to scale.\footnote{See also \cite{Ager2020a, Cermeno2019, Brakman2004Bombs}.} 
The empirical analysis here provides evidence from an abrupt change in first‑nature geography, permitting a causal estimate of these fundamentals. \cite{Bleakley2012} show that first‑nature effects persist. The findings presented below supplement that evidence: despite persistence, first‑nature geography continues to exert a meaningful effect. A literature review by \cite{Bosker2022} calls for natural experiments that isolate the impact of first‑nature geomorphology on long‑run economic outcomes—precisely the objective achieved in this paper.

% Contrasts with ancient‑economy studies; modern setting advantage
The natural experiment examined here is related to a few studies of ancient economies. \cite{Seror2020Random} uses random alterations in the course of the Yellow River to assess the impact of waterways on population growth in ancient China. \cite{Allen2023} employ river shifts and archaeological evidence from Mesopotamia to link first-nature geography to state formation. \cite{Matranga2024} links first-nature variation to the emergence of agriculture. Only historical settings such as these allow observation of the truly long-run effects of rare geographic shocks. The present paper is the first to do so in a context that also provides dense, high-quality microdata — a combination that simply does not exist elsewhere. Moreover, although early historical and Neolithic economies with simple settlements were understandably shaped by first-nature conditions, it is less clear whether the same holds with stronger second-nature geographical effects in a market economy \citep{Krugman1991, Krugman1999}. All of this makes the \textit{perfect storm} in the Limfjord both a rare and an ideal laboratory for precisely identifying the causal influence of first-nature geography on the location of economic activity.

% Identifies gap on geomorphology beyond seeds; claims contribution
Moreover, while previous research has robustly linked first‑nature features—such as soil quality and resource endowments—to economic outcomes through specific mechanisms,\footnote{\cite{Nunn2011a, HeavyPlough2016, Dalgaard2020, clark2014, ORourkeCoal2021, CoalPeatPaper}.} the role of first‑nature geomorphology remains underexplored outside studies of early civilizations.\footnote{\cite{Diamond1997, Allen2023, Matranga2024, Bakker2021Phonecians}.} This gap raises an important question: Does geomorphology matter beyond providing the initial seeds of future development? The evidence presented here answers this in the affirmative. 

I believe this paper provides one of the cleanest causal tests to date that geography’s effects persist beyond initial settlement patterns and continue to actively shape economic outcomes. In doing so, the geography-development debate must be shifted. Instead of asking how geography seeded divergence, we must ask how geography continues to influence both regional prosperity and disparities. 

% Roadmap
The remainder of the paper is organized as follows. The following section provides an overview of the historical background; Section 3 describes the data; Section 4 outlines the empirical strategy. Sections 5 and 6 present the main results on trade and population, respectively. Section 7 presents results for fertility and occupational structure. Section 8 examines the effect of the 12th‑century channel closure, and Section 9 concludes.

\begin{figure}[ht]
\centering
\caption{Map of Denmark and the Event in 1825}
\includegraphics[width=0.8\textwidth]{Plots/Map.png}
\parbox{0.9\textwidth}{
\caption*{\footnotesize \textit{Notes:} The map illustrates the effect of the 1825 Agger Isthmus breach on shipping routes. Improved market access is shown in the fill color, while the Limfjord regions are delineated by border color. The arrows indicate shipping routes before (dashed) and after (solid) the breach. \\ \textit{Source:} Parish borders from www.digdag.dk}
}
\label{fig:main_map}
\end{figure}


\FloatBarrier
\section{Historical background} \label{hist}

% Early‑19th‑c Denmark was poor; waterways (esp. Agger) key to later take‑off puzzle.
Denmark of the early 19th century was among the lowest income 
and least developed countries of Europe. Nowadays, Denmark is so well-developed that 'getting to Denmark' is the term used in the development literature for how to get to a place with modern well-functioning institutions and high living standards \citep{fukuyama2011origins}. The low-income status of Denmark would all change rapidly by the end of the century, and this sudden change is the subject of a growing literature.\footnote{\cite{Khaustova2015, Henriksen2006, milkandbutter, gorges2025tracksmodernity, BobergFazlic2023GettingtoDK}} According to \cite{Lampe2015DanesUK}, the connection with England was an important precursor for the later take-off.\footnote{As a piece of historical research this paper is a contribution to that question: How important was western market access for the Danish take-off?} Essentially this entire story is shaped by waterways, that allowed access to trade. But how and why was it important? And what specifically happened in the case of the Agger channel?

% Defines fjords; introduces Limfjord’s geography & Viking‑era importance; sets stage.
A fjord is a navigable inlet created by glacial activity during the last ice age, typically found in Northern Europe—particularly in Scandinavia—where they serve as natural highways connecting remote regions. The term “fjord” derives from Old Norse—and ultimately from the Proto-Indo-European root “pertu-” (also giving rise to “port” and the prefix “per-” as in per capita \citep{EtymFjord}) — underscoring fjords’ long-standing role in transportation. Indeed, many Scandinavian towns and cities—even today—are located along fjords.\footnote{Examples include Stockholm, Oslo, Gothenburg, Bergen, Stavanger, Odense, Aalborg, and Roskilde.} The present paper focuses on the Limfjord in Northwestern Denmark, located in the northern part of Jutland and positioned at the heart of Scandinavia between modern-day Norway and Sweden. During the Viking Age, the Limfjord featured both eastern and western openings that provided a safe shortcut for ships, facilitating trade and Viking expeditions. Unsurprisingly, it became a key gathering point for Vikings embarking on westward trade (and pillaging) \citep{Rasmussen1966}. In its center, the settlement of Aggersborg grew prominent as the site of a large ring castle built in the 10th century, symbolizing the consolidation of power by the emerging Danish state \citep{pedersen2014}. However, this period of prosperity would not last: The western opening eventually closed, and the local economy stagnated, as shown by the analysis in this paper.

% Dates closure of original west channel (1086‑1208) via historical evidence.
The last historical evidence of a western opening dates to 1085. King Canute IV gathered the last Viking fleet in the Limfjord to sail west and uphold his claim to the English throne, following in the footsteps of his great uncle, Canute II (“the great”), who once ruled a North Sea empire encompassing England, Denmark, Norway, and parts of Sweden \citep{Spejlborg2012}. However, his sailors and soldiers rebelled, the fleet never left Denmark, and Canute IV was killed while seeking refuge in Saint Alban’s Church in Odense in 1085 \citep{Pajung2012}.\footnote{The very city where I wrote this paper.} As a result, Danish Vikings never again attempted to conquer England \citep{Roesdahl2009}. Crucially for this study, this provides a range of dates for the closure of the previous channel: Between 1086 and 1208.\footnote{Further evidence for this can be found by cross-referencing Saxo \cite{saxo} [book XIII, section 5], \cite{Mortensen2018}, and geological evidence from \cite{Christensen2004}.}

% Shows economic shift: East Limfjord prospers, West stagnates due to lost access.
What followed was a complete reorientation of the Limfjord region’s prosperity. The Eastern Limfjord—particularly the town of Aalborg—emerged as a major center of the herring trade between Hanseatic Lübeck and the Baltic region \citep{sildeboom2022}. By 1672, Aalborg had become the largest Danish market town after Copenhagen, while towns in the Western Limfjord lagged far behind \citep{Degn1989}. Local institutions reflected these geographical constraints; the dangerous and costly approach to the west coast forced Western Limfjord market towns to access trade by traveling eastwards through the difficult waters of Løgstør and then onward to Aalborg (see Figure \ref{fig:main_map}). Several 16th- and 17th-century court cases reveal that these towns unsuccessfully fought for the right to trade independently of Aalborg merchants \citep[pp. 78-89]{ThistedLokalhistorie1974}, and overall, only limited trade occurred along the western routes \citep{Poulsen2019}.\footnote{Although some trade did occur via the perilous west coast under direct or indirect Aalborg authority—small boats risked the dangerous waters to sustain modest trade between Norway and the Western Limfjord; for example, in 1800 Thisted shipped 6,993 barrels of barley and 6,832 barrels of oats, with 31 percent and 47 percent respectively coming directly from the hazardous west coast beaches (the remainder via Aalborg). Most of this trade likely went to Norway, which in turn exported timber for construction \citep[p. 30, p. 234]{Aagard1802, Christensen1735}.}

% 1825 storm breach described; outlines three‑stage economic impact overview.
On the night between February 3rd and 4th, 1825, a storm flooded the narrow strip of land known as ‘Agger tange’ (Agger Isthmus) in a remote corner of the Limfjord. The sparsely populated Agger parish (with only 388 inhabitants in the 1801 census) saw its residents offered money to relocate to safer areas, and many settled further into the Limfjord \citep{Poulsen2019, Poulsen2022}. More importantly, the storm washed away the protective shrubs and plant material, fundamentally connecting the Limfjord to the North Sea by forming a shallow, navigable channel. The economic impact of the breach unfolded in three stages. \textit{First}, there was a windfall gain in fish catches (1825–1828) as fish, forced by the salinity change, became easier to catch—though fishing later stagnated in the Limfjord leading to North Sea fishing instead. \textit{Second}, once the channel became navigable in 1834, a large influx of trade ensued. \textit{Third}, this surge in trade spurred significant institutional and infrastructural adaptations, further boosting local prosperity from 1840 onward.

% First effect: salinity shock → short fishing boom then collapse/starvation.
As for the \textit{first effect}, the breach initially altered water salinity, forcing fish accustomed to brackish conditions into narrow rivers and streams where they were more easily caught. This led to a brief windfall in catches from 1825–1828—often called a “golden age” of fishing \citep{Poulsen2007} — after which yields fell sharply, several species declined or went nearly extinct \citep{Poulsen2019}. The local population suffered starvation for several years until the fishermen adapted to new opportunities in the North Sea. 

% Second effect: channel navigable 1834 → trade surge; stats on traffic & exports.
The \textit{second effect} was the increase in trade of the region. The channel’s navigability, achieved by 1834, triggered a substantial surge in trade. For instance, the market town of Thisted saw ship arrivals rise from 6 in 1834 to 62 in 1876, and exports increased dramatically—by 354 percent for barley and 997 percent for oats \citep[p. 153-159]{ThistedLokalhistorie1974}, compared to figures from 1800 \citep{Aagard1802}. Overall, ship traffic in the Limfjord grew from 19 vessels in 1835 to nearly 2000 in 1855 \citep{Svalgaard1977}, and newspaper reports of ships arriving from England further reaffirming this golden age of trade.\footnote{Example newspaper reports can be found in \citet{ThistedAmtsavis1834, RoskildeAmt1836, ViborgStift1852}}

% Third effect: trade boom spurs institutional & port investments (1840s‑1860s).
As for the \textit{third effect}, the influx of trade spurred significant institutional and infrastructural adaptations. As argued by \cite{Redding2015}, initial infrastructure investments often attract further improvements. The case of the 'natural infrastructure' of the Agger Channel studied here is also subject to this effect in many forms. In 1841, the King granted international trading rights to all Limfjord market towns—ending Aalborg's monopoly—and prompted the establishment of new customs offices and port facilities to manage the increased traffic.\footnote{For example, Thisted received a new port in 1841 \citep[p. 384-386]{Dioerup1842Thisted}; Struer’s port was expanded in 1856 and 1864 \citep[vol V, p. 467]{Trap3}; Lemvig built a port in 1857 \citep[vol V, p. 474]{Trap3}; Nykøbing Mors, with a port since 1788, was purchased in 1843 to support trade \citep[vol IV, p. 214]{Trap3}; and Løgstør’s port, originally from 1820, was expanded in 1852 \citep[vol IV, pp. 399-400]{Trap3}.}

% Further adaptations: canals, steamships, new port, protection works secure access.
The sudden channel formed by the 1825 storm also set off a cascade of other practical infrastructural adaptations, which are relevant to mention here. Fearing that the Agger channel would silt up, authorities first attempted to clear the shallow waters at Løgstør in 1843 \citep[p. 311, p. 4]{Bergsoee1844}—an effort that ultimately failed. Consequently, they commissioned the Frederik VII canal, completed in 1861 \citep{Petersen1877}, to maintain connectivity between the eastern and western Limfjord. Simultaneously, the new channel spurred the introduction of steamships in 1842, establishing the first ever Danish steamship route to England and further connecting Aalborg, Copenhagen, and other Limfjord towns,\footnote{\citet{Klem1967, Schovelin1891, Lampe2015DanesUK}} although the channel’s shallow depth limited vessel size \citep{Lassen1883}. As navigational challenges persisted and the Agger channel eventually silted up, Esbjerg emerged as a key port for Danish/British trade \citep{Lampe2015DanesUK}. Then, following another storm in 1862, a new channel—the Thyborøn channel—opened in 1867 in approximately the same location and remains in use today \citep{Petersen1877, Ravn1993}. In response to mounting political pressure and fears of further siltation, additional coastal protection measures such as groins were installed along the northwestern coast in 1875 \citep{petersen1853oplysende, Trap3}. Today, Thyborøn, located at the channel's mouth, is one of Denmark's main fishing ports, accounting for 25 percent of the Danish catch in 2021 \citep{MinisterietforFodevarer2022}. And throughout the 20th century a strong relatively strong textile industry prospered in the region \citep{MirakletPaaHeden2025}. The origins of both can be seen emerging in the analysis of this paper. 

Together, the direct effects as well as the sequential adaptations — directly triggered by the sudden alteration of first nature geography in 1825 — secured long-term market access and laid the foundation for the economic transformation examined in this paper.

\section{Data}

% Describes Sound Toll Registers & channel‐traffic series used to measure trade (1750–1855).
For trade data, the analysis relies on the Sound Toll Registers from Elsinore, north of Copenhagen, which recorded detailed accounts of all ships taxed from 1420 to 1857—amounting to about 1.8 million entries \citep{gobel2010oresundstolden, soundtoll_data}. These records include information on the origin and destination of each ship, with ports conveniently labeled with geographic coordinates. Based on this, the ships traveling to or from 126 ports in present-day Denmark are counted, and this in turn is used as a measure of levels of trade. The extract covers the period 1750–1855. This dataset is further supplemented by archival data on the number of ships passing the newly formed channel from 1834 \citep{Svalgaard1977}.

% Outlines census data (Link Lives), parish crosswalk, HISCO coding and data quality checks.
Census data were obtained from Link Lives \citep{mathiesen2022linklives} and include individual-level information on occupation, age, gender, and parish of residence. For the analysis, these data are aggregated into parish-level population counts for the years 1787, 1801, 1834, 1840, 1845, 1860, 1880, and 1901. I manually constructed a crosswalk between the census parish names and the historical parish borders (available in the project's public repository), retaining only the 1589 parishes that appear in all censuses (out of 1783). In addition to population counts, the census provides occupational descriptions that have been standardized into HISCO codes using \textit{OccCANINE} \citep{dahl2024breaking}.\footnote{The full micro-level HISCO-coded census data are publicly available from \citep{dk_hisco_data}}. 200 random HISCO codes from the algorithm were manually and diligently checked to be 94 percent accurate.

% Explains summary‑stat variables, table & balancing‑plot; notes pre‑event differences handled 
Table \ref{tab:desc_pop} presents summary statistics for the main variables used in the analysis. “Population” denotes the number of people in each parish for a given census year. “HISCO Agricultural” counts individuals with HISCO codes beginning with 6, while “HISCO Manufacturing” counts those with HISCO codes starting with 7, 8, or 9. “Born in different county” records the number of residents born outside the county in which they later lived, and “Child–women ratio” (an indicator of fertility) is defined as the ratio of children aged 1–5 to women aged 15–45. Figure \ref{fig:bal} displays the standardized distributions of these variables before the Agger channel breach (with the exception of “Born in different county,” which is only available from 1845). Note that while there is substantial overlap between the distributions before the event, pre-event differences are also evident; these differences are addressed in Section 7 and further in Appendix \ref{dr_estimates}.

% Introduces archaeological dataset, geo‑matching to parishes, and panel construction.
Archaeological data were obtained from a public registry of Danish archaeological findings. This database contains geo-referenced records of sites that are categorized by type and dated to specific intervals. Each site's coordinates were matched to a parish using the parish borders described above, and the dating information was used to construct a panel of economic activity. Details of this are outlined in the relevant Section \ref{arch}.\looseness=-1

\FloatBarrier

\begin{table}
\centering
\caption{Summary statistics for parish level census data}
\label{tab:desc_pop} 
\footnotesize
\begin{tabular}{lcccccccc}
\toprule
  & Observations & Mean & SD & Min & Median & Max\\
\midrule
Population & 14301 & 652.02 & 505.88 & 27.00 & 526.00 & 13087.00\\
Affected: West Limfjord & 14301 & 0.12 & 0.33 & 0.00  & 0.00 & 1.00\\
Affected: $\Delta log(MA_i)$ & 14301 & 0.07 & 0.06 & 0.01 & 0.05& 0.27\\
HISCO Agricultural & 14301 & 181.07 & 120.03 & 0.00 & 152.00 & 1512.00\\
HISCO Manufacturing & 14301 & 65.65 & 74.77 & 0.00 & 46.00 & 2809.00\\
Born in different county & 7945 & 253.13 & 378.91 & 0.00 & 72.00  & 4299.00\\
Child-women ratio & 14274 & 0.47 & 0.11 & 0.00 & 0.46 & 1.36\\
\bottomrule
\end{tabular}
\parbox{0.9\textwidth}{
\caption*{\footnotesize \textit{Notes:} This table contains summary statistics for the variables used, which are ultimately sourced from the census data. For most variables, there are 14,301 observations corresponding to 1,589 parishes and 9 census years. The child–women ratio is defined as the ratio of children aged 1–5 to women aged 15–45; 27 observations are missing from this variable since no women in the relevant age group are observed. 'Born in different county' is only observed from 1845 onward. \\ \textit{Source: Danish census data}}
}
\end{table}

\begin{figure}
\centering
\caption{Variable distributions}
\label{fig:bal}
\includegraphics[width=1\textwidth]{Plots/Balancing_plot.png}
\parbox{0.9\textwidth}{
\caption*{\footnotesize \textit{Notes:} This figure shows the distribution of variables of interest before the 1825 breach in the West Limfjord and the rest of the country (excluding other Limfjord parishes). Overall, the distributions overlap substantially, though the West Limfjord generally exhibits lower population density and fertility. \\ \textit{Source: Danish census data}}
}
\end{figure}

\FloatBarrier

\section{Empirical Strategy} \label{emp_strat}

% Introduces natural experiment, lists outcome dimensions, states TWFE DiD model.
This paper exploits a natural experiment generated by the 1825 Agger channel breach to identify the causal impact of an abrupt change in first–nature geography on economic outcomes. Several dimensions are considered: the effect on trade, the effect on population dynamics—including changes in fertility, migration, and occupational structure—and the external validation provided by a reverse natural experiment from the medieval period. A difference–in–differences design is employed to compare outcomes in areas directly affected by the channel to those that were not. The empirical strategy, which is similar across all outcomes, follows a standard TWFE specification:\footnote{Note that despite recent debate about TWFE and DID the specific setting of this paper is such that the simple TWFE is appropriate \citep{roth2023trendingDiD}.}
\begin{equation}
\label{eq:eq501}
\log(y_{it})= \gamma_t + \gamma_i + \textit{Affected}_{i}\,\beta_{t} \,  + \varepsilon_{it},
\end{equation}

\noindent
where, $y_{it}$ is the outcome of interest (e.g., population) for parish $i$ at time $t$. $Affected_{i}$ measures the degree to which parish $i$ was affected by the new channel. $\beta_t$, the parameter of interest, captures the year specific effect. $\gamma_t$ and $\gamma_i$ are year and parish fixed effects. Generally, the last observed year before the event is used as the reference year.  

% Defines first treatment measure: West‑Limfjord dummy via geographic heurestic
$\textit{Affected}$ is measured using two approaches. First, a simple dummy measure categorizes parishes based on their geographic location relative to the Limfjord. Parishes are classified as “Limfjord” if they are closer to the Limfjord than to any other coast. Among these, a roughly northwest–southeast line\footnote{Defined by the coordinates [57.044185, 9.186837] and [56.958951, 9.275585].} separates the West from the East, with a 20 km buffer zone designating the Middle Limfjord. The area, which was principally affected was the West Limfjord. As such "Affected" is defined as being in the West Limfjord, while the Middle and East Limfjord are also added as controls in the dummy specification. 

% Defines second treatment measure: continuous change in market access via cost–distance formula.
Second, \textit{Affected} is defined from a simple measure of the change in market access implied by the geographical change. Market access for each parish is computed using a cost–distance approach that reflects the shift in the set of effective ports from before to after the isthmus breach. In this framework, market access is given by

\begin{equation}
\label{eq:ma_eq}
\textit{MA}_p = \sum_{h \in \mathscr{H}} \left[CostDist(p, h; \alpha) + 1\right]^\theta,
\end{equation}

\noindent
where $\mathscr{H}$ is the set of ports available, $CostDist(p, h; \alpha)$ is the relative cost of traveling from parish $p$ to port $h$ (with over-land travel assumed to cost $\alpha = 10$ times as much as sea travel), and $\theta$ is the distance elasticity (assumed to be $-1$, following \cite{Harris1954, rauch2022a}). The change in market access is measured as $\Delta \log(\textit{MA}_p) = \log(\textit{MA}_{p,\textit{after}}) - \log(\textit{MA}_{p,\textit{before}})$, where $\mathscr H$ changes between \textit{before} and \textit{after} to include the newly available ports, capturing how the breach transformed the regional network of ports. For additional technical details, see Appendix \ref{details_ma}, and for robustness using different baseline parameters see Appendix \ref{pop_multiverse}. Figure \ref{fig:main_map} in Section \ref{intro} illustrates the change in market access across Denmark given the newly available ports.



\FloatBarrier
\section{The Effect on Trade} \label{trade}

% Descriptive figures show sharp post‑1834 surge in West/Middle Limfjord traffic despite Napoleonic dip.
The opening of the Agger channel in 1825 fundamentally transformed regional trade in the Limfjord area — especially in the West and Middle Limfjord. Figure \ref{fig:Sound_toll} shows the log-transformed total traffic (inbound and outbound) recorded by the Sound Toll Registers. Although these registers capture only ships traveling eastward through the toll at Elsinore, the data reveal that ports in the West Limfjord experienced very low traffic before 1825, followed by a marked and sustained increase after the channel opened - especially after 1834 when the channel became fully navigable (see Section \ref{hist}).  This pattern persists despite a temporary disruption during the Napoleonic Wars (1807–1814) \citep{Feldbaek2015}. Figure \ref{fig:channel} reinforces this observation by presenting the actual number of ships passing through the Agger channel (and later, the Thyborøn channel). The takeaway from this is that the new channel caused a large influx of trade-activity to the affected region. 


\begin{figure}
\begin{center}
  \caption{Number of ships - sum of inbound/outbound)} \label{fig:Sound_toll}
  \includegraphics[width=1\textwidth]{Plots/Ship_trafic.png}
  \parbox{1\textwidth}{
  \caption*{\footnotesize \textit{Notes:} This shows the log-transformed sum of traffic to and from ports in Denmark as captured by those ships that passed Elsinore. \\ \textit{Source: Sound Toll Registers.}}
}
\end{center}
\end{figure}

\begin{figure}[ht]
\begin{center}
  \caption{Number of ships passing the Agger channel} \label{fig:channel}
  \includegraphics[width=1\textwidth]{Plots/Ship_trafic_channel.png}
  \parbox{1\textwidth}{
    \caption*{\footnotesize \textit{Notes:} This shows the number of ships passing the Agger channel and later Thyborøn channel. \\ \textit{Source: Svalgaard (1977). The single observation from 1867 is from Ravn (1993).}}
  }  
\end{center}
\end{figure}

% Outlines DiD/PPML specification, treatment dummy, robustness checks for trade effect.
To formally test the effect of the channel on trade, the simple dummy measure of being affected (defined in Section \ref{emp_strat}) is used in a difference–in–differences framework. In this specification, the amount of trade at port $i$ in year $t$ is modeled as a function of time fixed effects, port fixed effects, and an interaction between a post-1834 indicator and the affected dummy variable. The estimation is conducted using the Poisson Pseudo Maximum Likelihood (PPML) estimator, which has become the standard for this type of regression which inherently has many zeros in the outcome \citep{Silva2006, Marczinek2022, SantosSilva2022}. To address potential distortions from the Napoleonic Wars and the early transition period (1825–1833), results are also reported where these years are excluded. Additional robustness checks include estimates using alternative transformations ($\log(y+1)$ and the inverse hyperbolic sine) as well as an analysis of the extensive margin (i.e., whether any trade occurred).
\FloatBarrier

\begin{table}
\centering
\caption{Channel Introduction and Trade} \label{tab:reg_trade}
\footnotesize
\begin{tabular}{lccccc}
   \tabularnewline \midrule \midrule
   Dependent Variables: & \multicolumn{2}{c}{traffic}     & log(traffic+1)  & arcsinh(traffic)  & $1[traffic>0]$\\
                        & (1)            & (2)            & (3)             & (4)             & (5)\\  
                        & Poisson        & Poisson        & OLS             & OLS             & OLS\\  
   \midrule
   Post$\times$east     & -0.4896$^{**}$ & -0.4817$^{*}$  & -0.0029         & 0.0115          & 0.0115\\   
                        & (0.2172)       & (0.2489)       & (0.0915)        & (0.1741)        & (0.1741)\\   
   Post$\times$middle   & 5.557$^{***}$  & 11.76$^{***}$  & 0.9672$^{***}$  & 0.3811$^{***}$  & 0.3811$^{***}$\\   
                        & (0.2151)       & (0.2471)       & (0.0584)        & (0.0233)        & (0.0233)\\   
   Post$\times$west     & 4.129$^{***}$  & 4.698$^{***}$  & 0.3821          & 0.3194$^{**}$   & 0.3194$^{**}$\\   
                        & (0.5090)       & (0.4714)       & (0.2635)        & (0.1263)        & (0.1263)\\   
   \midrule
   Observations         & 13,356         & 11,214         & 13,356          & 13,356          & 13,356\\  
   \midrule
   1807-1814 excl.      & No             & Yes            & No              & No              & No\\
   1825-1833 excl.      & No             & Yes            & No              & No              & No\\
   \midrule \midrule
\end{tabular}
\parbox{0.95\textwidth}{
\caption*{\footnotesize \textit{Notes:} Cluster-robust standard errors (clustered at the port level) are reported in parentheses. Column (1) shows Poisson estimates; Column (2) excludes years 1807–1814 and 1825–1833; Columns (3) and (4) use OLS with $\log(y+1)$ and arcsinh transformations to address zero outcomes; Column (5) analyzes the extensive margin (whether any trade occurred). \\ \textit{Source: Sound Toll Registers Online}}
}
\end{table}

% Reports regression table: large, robust trade increase supports descriptive pattern and narrative.
Table \ref{tab:reg_trade} summarizes the regression estimates. Across all specifications, the results consistently indicate what is also shown in the descriptive statistics: Ports in the West and Middle Limfjord experienced a large increase in trade after 1834. The estimated coefficients remain robust when excluding data from disruptive periods, and alternative estimators yield qualitatively similar results. These findings support the historical narrative that the opening of the Agger channel dramatically increased regional trade, with the resulting surge in maritime traffic placing additional pressure on local infrastructure and spurring further investment (See Section \ref{hist}). In sum, first–nature geography played a critical role in shaping regional trade patterns.

\FloatBarrier
\section{The Effect on Population Density}

\begin{figure}
    \centering
    \caption{Effect of the Agger channel on population size}
    \begin{subfigure}[b]{0.8\textwidth}
        \centering
        \caption{Dummy approach}
        \includegraphics[width=0.7\textwidth]{Plots/Regression_plots/pop_dummy.png}
    \end{subfigure}
    \vspace{1cm}
    \begin{subfigure}[b]{0.8\textwidth}
        \centering
        \caption{Market access approach}
        \includegraphics[width=0.7\textwidth]{Plots/Regression_plots/pop_MA.png}
    \end{subfigure}
    \parbox{1\textwidth}{
    \caption*{\footnotesize \textit{Notes:} Effect of the Agger channel on population size in affected areas. The error bars represent 95\% confidence intervals based on cluster-robust standard errors. Panel (a) uses a dummy measure (being in the West Limfjord), and panel (b) uses a measure based on the estimated improvement in market access. The year 1801 is the reference. All parameter estimates are reported in Appendix \ref{pop_reg_table}. \\ \textit{Source: Danish census data}}
    }
    \label{fig:pop1}
\end{figure}


% Regression results: By 1901 pop ↑26.7 % (dummy) and elasticity 1.59 (MA).
The improved first–nature geography brought about by the Agger channel breach led to a significant increase in population density. Figure \ref{fig:pop1} presents regression results estimating this effect. In panel (a), where being affected is defined as a dummy for being in the West Limfjord, the results indicate that by 1901 the first-nature geographical shock had caused a 0.236 log-point (26.7 percent) larger populations in the affected parishes. Panel (b) shows that when "Affected" is defined by the estimated improvement in market access, the effect corresponds to a market access elasticity of 1.59 by 1901 — comparable to the magnitude found for the Panama Canal \citep{rauch2022a}.

% Effect is delayed until 1860; aligns with three‑stage historical mechanism.
It is noteworthy that the population effect was not immediate. Despite a sharp increase in ship traffic following the channel opening (Section \ref{trade}), a positive effect on population size only becomes discernible from 1860 onward. In fact, using the market access approach, there is a slight negative effect immediately after the channel became navigable — consistent with the salinity shock and associated food shortages described in the historical background. Thus, the delayed demographic response is unsurprising and closely mirrors the three-stage adjustment process outlined in Section \ref{hist}. More broadly, this result illustrates that first-nature geography operates on a long-run time scale: even when underlying geographic conditions shift abruptly, their full economic consequences may take decades to materialize. In this case, the effect unfolds over 76 years, from 1825 to 1901.

% Robustness: multiverse shows no pre‑trends; doubly‑robust DiD confirms balance.
Robustness of these findings is tested against a series of alternative specifications. Figure \ref{fig:pop2} in Appendix \ref{pop_multiverse} presents the multiverse of results based on different comparison groups and parameter choices. None of the alternative specification show any pretrends. Another potential worry is balance. This is addressed by correcting for covariates (age, occupation, fertility) using the doubly-robust difference-in-differences estimator suggested by \cite{Callaway2021did}. The results of this are shown in Appendix \ref{dr_estimates}. The effect of the channel is robust to all of this.

% Synthesizes: strong evidence that channel increased population density.
Together, these results provide robust evidence that the Agger channel breach led to a significant increase in population density in the affected areas.

\section{Fertility, Migration and Occupations} \label{further_res}

% Motivates section: test if pop growth tied to prosperity; claims manufacturing & fishing ↑; fertility, not migration, drove growth.
For population growth to occur, the people must come from somewhere. This section tests whether the observed population growth plausibly relates to higher levels of prosperity. This is done to corroborate the connection between the shock to first-nature geography and economic development. It is demonstrated that both the manufacturing and fishing industries grew, indicating that the West Limfjord fundamentally became more prosperous. Furthermore, consistent with the post-Malthusian regime of the time \citep{Jensen2022, Klemp2016}, evidence is shown that population growth is explained by fertility rather than internal migration. This also indicates that the population growth was intrinsic rather than reallocated from elsewhere.

\subsection{Occupations}

% Sets up occupation analysis: HISCO codes, DiD framework, zero-outcome challenge.
To assess the impact of the channel on occupational structure, the HISCO coding system \citep{leeuwen2002hisco} is employed, in which the first digit of each code denotes a major occupational category. The effect on each category is estimated using a difference–in–differences framework similar to that presented in Equation \ref{eq:eq501} (Section \ref{emp_strat}). However, in this setting there are frequent zero-outcomes. Zero outcomes imply a methodological trade-off between measuring the intensive and extensive margins separately or accepting that the measured coefficients are not unit-invariant \citep{roth2023loglike}. The functions $f_k()$ addresses this issue in four distinct ways:
\begin{enumerate}
    \item Extensive: $f_1(y_{it}) = 1[y_{it}>0]$
    \item Intensive: $f_2(y_{it}) = \log(y_{it})$ (restricted to observations with $y_{it}>0$)
    \item Combined (log-transform): $f_3(y_{it}) = \log(y_{it}+1)$
    \item Combined (arcsinh-transform): $f_4(y_{it}) = \operatorname{arcsinh}(y_{it})$
\end{enumerate}

% Introduces four transformations f_k for extensive/intensive/combined margins.
The extensive margin captures the proportion of parishes with at least one individual in a given occupation, while the intensive margin reflects the magnitude of the effect among parishes where the occupation is already present. The third and fourth transformations provide standard approaches to combining these margins, acknowledging that the estimates are not unit invariant.

% Defines APE_share metric to scale effects relative to parish population.
A key concern is whether the observed changes are large relative to the population size. Interpreting the relative effects is not sufficient if, for example, the number of priests increases by 50 percent from two to three, then that might not meaningfully be described as a shift of the occupational structure. Therefore, all estimates are converted into an \textit{Average Partial Effect share} ($APE\,share$), which represents the average number of individuals gained in an occupation in an average parish, expressed as a share of the average West Limfjord parish population in 1901. This is computed using the formula:

\begin{equation}
    \label{eq:ape_share}
    APE\,share_j = \frac{\overline{Occ_{j, 1901}}}{\overline{Pop_{1901}}} \hat{\beta_j},
\end{equation}

\noindent
where $\overline{Occ_{j, 1901}}$ is the average count of individuals in occupation $j$, $\overline{Pop_{1901}}$ is the average parish population in the West Limfjord in 1901, and $\hat{\beta_j}$ is the estimated effect for occupation $j$.

% Explains multiple-testing issue; limits focus to 56 estimates; Bonferroni correction.
Given that testing all seven occupational categories across the dummy and market access measures, as well as four specifications of $f_k()$, across multiple census years results in a large number of parameters (504 in total), this presents a problem of multiple testing. Therefore, only the 56 estimates for 1901 are presented. To account further for multiple testing, the standard errors of the remaining coefficients are adjusted using the conservative Bonferroni correction.

% First figure results: intensive increase in Agriculture & Manufacturing; population growth linked to fertility.
Figure \ref{fig:mech_occ} displays the Bonferroni-corrected results using the dummy approach, showing the estimated effect on each major occupational category in 1901. The plot reveals a significant impact on the intensive margin for two categories: Agriculture (HISCO codes starting with 6) and Manufacturing (codes starting with 7, 8, or 9). These findings suggest that, alongside overall population growth driven by increased fertility (as discussed in the following subsection), there was a substantial absolute shift in the occupational structure, with a marked rise in manufacturing — a sector that was previously uncommon. Detailed coefficients from both the dummy and market access approaches, as well as event study plots, are available in Appendix \ref{all_occ_results} and the project's online repository.\footnote{\url{https://github.com/christianvedels/A_perfect_storm/tree/main/Plots/Mechanism/Occupations}}

% Second figure detailed breakdown: industrialization occupations (textile, generic factory) ↑ ; fishing also ↑ .
Figure \ref{fig:mech_occ2} provides a more detailed breakdown by using the first three digits of the HISCO codes for agricultural occupations and the first two digits for manufacturing occupations. Due to the large number of regressions conducted, confidence intervals for Figure \ref{fig:mech_occ2} were not estimated, so the figures should be interpreted as descriptive and indicative. Notably, manufacturing occupations that typically signal industrialization — such as textile production and generic factory work ("workers not elsewhere classified") — show a large growth, indicating a marked increase in economic activity driven by the improved first–nature geography. In the agricultural sector, the rise in fishing occupations reflects the benefits of the sudden availability of fish from the North Sea. 

% Synthesis
Overall, these findings provide evidence that the improved first–nature geography contributed to a shift in the occupational structure associated with higher levels of productivity.

\begin{figure}
\begin{center}
  \caption{Impact of the Agger Channel on Occupational Structure in 1901} \label{fig:mech_occ}
  \includegraphics[width=0.8\textwidth]{Plots/Mechanism/All_occupations_dummy.png}
  \parbox{1\textwidth}{
  \caption*{\footnotesize \textit{Notes:} This figure illustrates the average partial effect (as a share of parish size) on the occupational structure in 1901. The first panel shows the effects on the extensive margin. The second panel shows the effect on the intensive margin (with included number of parishes shown as well). Instead of choosing between intensive and extensive margins, the third panel uses the $log(x+1)$ transformation and the fourth panel uses the inverse hyperbolic sine transformation. The results are based on the Dummy definition of being affected by the channel. Appendix D.1 shows a table of all results including results based on market access. The error bars represent 95 percent confidence intervals corrected for multiple testing using the Bonferroni correction. This is based on standard errors, which are clustered at the parish level.}
}
\end{center}
\end{figure}


\begin{figure}
\begin{center}
  \caption{Effects on Detailed Occupational Structure} \label{fig:mech_occ2}
  \includegraphics[width=1\textwidth]{Plots/Mechanism/Detailed6789/Dummy_asi.png}
  \parbox{1\textwidth}{
  \caption*{\footnotesize \textit{Notes:} This figure depicts the relative changes in subcategories of agricultural and manufacturing occupations (HISCO codes starting with 6 and 7/8/9). The plot presents the results using the arcsinh transformation and the dummy-definition of being affected by the channel. Qualitatively similar results using alternative approaches are available in the online repository.}
}
\end{center}
\end{figure}

%\FloatBarrier
\subsection{Fertility and Migration}

% Introduces Fertility vs Migration question and theoretical expectation (post‑Malthusian).
A remaining question is whether the observed population growth resulted from fertility or migration. This in turn, given Malthusian dynamics, is indicative as to whether the population growth came from an intrinsic improvement in prosperity or from the reallocation of prosperity from elsewhere. If the channel merely attracted people from elsewhere, increased internal migration would be expected in the affected areas. In contrast, post-Malthusian dynamics \citep{Jensen2022, Klemp2016} predict that improved living standards lead to higher fertility.

% Event‑study results: fertility ↑ 11.6 %, migration flat/↓; elasticity ~0.96.
Figure \ref{fig:migr_fert} displays event-study plots for both fertility and migration. Fertility is measured by the child–women ratio, while migration is defined as the share of residents born in a different county than their current county. Panels (a) and (b) show that by 1901, the channel was associated with an 11.6 percent increase in fertility —corresponding to an elasticity of approximately 0.96 percent per one percent increase in market access. In contrast, panels (c) and (d) indicate that migration either declined or remained unchanged, suggesting that the population increase was not driven by reallocation. 

% Explains timing pattern and age‑composition effect; concludes growth intrinsic via fertility.
The evolution of the fertility effect mirrors the historical background: The initial surge in fish catches following the channel breach boosted fertility as early as 1834, followed by a collapse in within-Limfjord fishing and subsequent adaptation to open-ocean fishing and enhanced trade opportunities. Detailed results by age group are provided in Appendix \ref{effects_by_age_group}. These results show that the population composition became comparatively younger as more children were born and survived beyond infancy. All of this indicates that population growth came from an intrinsic improvement in the prosperity of the region which was affected. 

\begin{figure}
    \centering
    \caption{Effects on fertility and internal migration}
    \begin{subfigure}[b]{0.45\textwidth}
        \centering
        \caption{Fertility (MA approach)} \label{fig:fert_ma}
        \includegraphics[width=\textwidth]{Plots/Mechanism/fertility_MA.png}
    \end{subfigure}
    \hfill
    \begin{subfigure}[b]{0.45\textwidth}
        \centering
        \caption{Fertility (dummy approach)} \label{fig:fert_dummy}
        \includegraphics[width=\textwidth]{Plots/Mechanism/fertility_Dummy.png}
    \end{subfigure}
    \vspace{0.45cm}
    \begin{subfigure}[b]{0.45\textwidth}
        \centering
        \caption{Internal migration (MA approach)} \label{fig:migr_ma}
        \includegraphics[width=\textwidth]{Plots/Mechanism/born_different_share_MA.png}
    \end{subfigure}
    \hfill
    \begin{subfigure}[b]{0.45\textwidth}
        \centering
        \caption{Internal migration (dummy approach)} \label{fig:migr_dummy}
        \includegraphics[width=\textwidth]{Plots/Mechanism/born_different_share_Dummy.png}
    \end{subfigure}
    \parbox{0.9\textwidth}{
    \caption*{\footnotesize \textit{Notes:} This shows the effect on indicators of fertility and migration. Panel (a) and panel (b) show the effect on the child-women ratio. Panels (c) and (d) show the effect on the number of people born in a different county, than where they usually live as a share of the total population in that parish. The results of panels (a) and (b) indicate that the channel caused fertility to increase. The results of panels (c) and (d) indicate, that migration either contributed negatively to the increase in population or did not contribute at all.  \\ \textit{Source: Danish census data}}
}
    \label{fig:migr_fert}
\end{figure}






\FloatBarrier

\section{The Reverse Natural Experiment} \label{arch}

\subsection{Method and data}
% Sets stage: medieval channel closure as reverse natural experiment; external‑validity motivation.
The Limfjord region experienced a favorable first–nature geography after 1825, but this was not the first time the region had benefited from a natural channel. Prior to the 1100s, a similar channel existed in nearly the same location (see Section \ref{hist}), only to close sometime between 1086 and 1208 — in effect, the same natural experiment occurred, but in reverse. Medieval Denmark, freshly converted from paganism and dominated by Viking traditions without modern scientific advances, contrasts sharply with 19th–century Denmark in every respect except for the channel’s location. In this way, comparing these two periods is akin to attempting replication across two entirely different countries, while the only constant is the first-nature geography and the presence (or absence) of the channel. As such this offers a natural way of testing external validity: Will the same effect replicate in such a (culturally, technologically, institutionally) different environment?

% Justifies using archaeological finds as economic proxy via Settlement Scaling Theory.
A significant challenge in studying this medieval shock is the sparsity of available data. To overcome this, archaeological evidence is used as a proxy for economic activity or population density. This approach draws on Settlement Scaling Theory, which posits that larger settlements yield more archaeological finds \citep{Ortman2020}, and this approach has also seen recent application in the economics literature.\footnote{\cite{Davis2002, Bakker2021Phonecians, Allen2023, Barjamovic2019}} 

% Describes archaeological registry, filters to 750‑1500 CE coins & buildings, yields sample sizes.
The archaeological data come from the Danish national registry 'Fund og Fortidsminder', which reports all sites documented by Danish museums and the National Agency of Culture and Palaces. Each entry includes a geographic coordinate, a site type (e.g., coin, structure, burial mound), and a date range. The full dataset includes 290,524 findings, spanning from the Paleolithic to the modern period. For the purposes of this study, the focus is narrowed to entries dated between 750 and 1500 CE and restricted to two types most relevant to economic life: coins and buildings. Coin findings are interpreted as evidence of trade, while buildings proxy for population density. This filtered dataset yields 3,411 coin findings and 4,396 building findings, which are used to construct a panel of estimated parish-level economic activity.

% Explains naive panel construction bias with wide date ranges using example.
To clarify the construction of this panel, consider an example: Finding ID 338 is a coin discovered within the historic parish of Pedersker on the island of Bornholm, dated between 1300 and 1535. This suggests that economic activity involving coins occurred in that area during this interval. A simple method would record a value of '1' for Pedersker in each year between 1300 and 1535. Repeating this for all findings would yield a time-varying indicator of archaeological activity across parishes. However, this approach introduces bias: findings dated with broader year ranges would contribute to more years, and thus be overrepresented, even though their temporal precision is lower.

% Introduces probabilistic dating + Monte Carlo & bootstrap method; links to DiD framework.
Instead, the dating ranges are interpreted as probability distributions.\footnote{This dating range can reasonable be interpreted as a uniform distribution (which is the basis of the results shown here) or a normal distribution (for which results are shown in Appendix \ref{norm_arch}).} The probability that a finding corresponds to a particular year is then used to estimate the likelihood that any archaeological activity occurred in a given parish-year. The analytical derivation of this method — using Bayes’ rule to estimate the inverse probability $P(\text{finding}|t)$ — is detailed in Appendix \ref{math_note}. But it consists of following the logic of the archaeological dating uncertainty and using a simple Monte Carlo method to sample from the probability distribution, which describes the probability of a finding-generating activity occurred in a specific year. The end result is a panel of the probability economic activity across parishes and years, which enters the difference-in-differences framework outlined in Section \ref{emp_strat}. Inference is made with a clustered bootstrap, which is modified to take into account the sampling uncertainty of the procedure.\footnote{In brief, each iteration resamples (clusterwise) in the Monte Carlo samples that was used in panel construction.}

% Raises soil‑type confound; proposes propensity‑score matching with XGBoost on soil categories.
A potential concern is that soil types might affect both economic development and the preservation of archaeological remains. For example, the western Limfjord is known for its fertile soils, which may have led to distinct developmental paths \citep{HeavyPlough2016, WinnersAndLosers2022}, while soil conditions could also influence the survival of archaeological sites. To address this, propensity score matching is employed to compare parishes affected by the channel closure with those not affected but similar in soil type. Propensity scores are estimated using extreme gradient boosting trees \citep{chen2016xgboost}, a machine learning technique that generalizes random forests \citep{Breiman2001}, using soil type data from \citep{Pedersen2019}. Only soil types present in at least 10 percent of parishes are included to avoid overfitting on rare characteristics. 

% Details greedy matching and shows almost identical prop‑score distributions (Fig. prop).
Each West Limfjord parish is then matched with a counterpart outside the Limfjord region with the most similar propensity score, using a greedy, without-replacement algorithm that ensures each affected parish is uniquely paired with one control parish. Figure \ref{fig:prop_score} shows the propensity score distributions before and after matching, demonstrating that the matched groups are nearly identical in terms of soil characteristics.

% Presents descriptive coin‑count trends: West stagnates post‑closure, East/Denmark rise.
However, before examining the more intricate results, it is relevant to establish a descriptive pattern based on simple counts of findings. Figure \ref{fig:arch_desc} displays the average number of coin findings by year across the West, Middle, and East Limfjord regions as well as the rest of Denmark, without accounting for dating range uncertainty. The plot reveals that, following the closure of the channel between 1086 and 1208, coin findings gradually increased in the East Limfjord and the rest of the country, while remaining consistently lower in the West and Middle Limfjord. This pattern suggests that the West and Middle Limfjord experienced an economic and commercial stagnation during this period.

\begin{figure}
    \centering
    \caption{Soil type propensity scores}
    \begin{subfigure}[b]{0.45\textwidth}
        \centering
        \caption{Propensity score before matching} \label{fig:prop1}
        \includegraphics[width=\textwidth]{Plots/Propensity_before.png}
    \end{subfigure}
    \hfill
    \begin{subfigure}[b]{0.45\textwidth}
        \centering
        \caption{Propensity score after matching} \label{fig:prop2}
        \includegraphics[width=\textwidth]{Plots/Propensity_after.png}
    \end{subfigure}
    \parbox{0.9\textwidth}{
    \caption*{\footnotesize \textit{Notes:} Propensity scores between the West Limfjord and the reference group before and after matching. Propensity scores are estimated using extreme gradient boost. Matching is done with a \textit{greedy} matching procedure in random order.}
}
    \label{fig:prop_score}
\end{figure}


\begin{figure}
    \centering
    \caption{Rate of coin findings}
    \label{fig:arch_desc}
    \includegraphics[width=\textwidth]{Plots/Arch_descriptive.png}
    \parbox{0.9\textwidth}{
        \caption*{\footnotesize \textit{Notes:} This plot shows the average number of coin findings per year, counted without adjusting for dating range uncertainty, for the West, Middle, and East Limfjord regions, as well as for the rest of Denmark. \\ \textit{Source: Danish registry of archaeological findings}}
    }
\end{figure}

\FloatBarrier
\subsection{Results}
% Introduces regression outputs (fig/table) and outcome definition (±25 yr probability).
Figure \ref{fig:arch_reg} and table \ref{tab:arch1} provide estimates of the effect to the probability of archaeological findings. Figure \ref{fig:arch_reg_boot} shows the parameter distribution in 1350 yielded by the sampling procedure presented above and in Appendix \ref{math_note}. The regression table presents the results of archaeological regressions for the full sample of all of Denmark and a matched sample, using both the dummy definition of being affected and the change in market access approach. The outcome is the probability that a coin finding or building was generated in the area covered by that parish within $\pm 25$ years of the reported year. 

% Quantifies market‑access semi‑elasticities and dummy effects for 1350; coins & buildings.
The parameters from the Market Access approach should be interpreted as semi-elasticities, i.e. how does the probability of a coin being generated change with a relative change in market access? A 1 percent decrease in Market Access led to a 13-19 percent change in the probability of a coin finding being generated in 1350 - 150 years after the channel closed (columns 1 and 5). Equivalently, it led to a reduction in the chance of buildings being generated as archaeological findings by 6 to 8 percent (columns 3 and 7). The parameters for the dummy approach can be interpreted as the change in the probability that findings were generated from a West Limfjord parish compared to other parishes. In 1350, the chance of coin findings being generated had fallen by around 1.5 or 1.9 percent  (columns 2 and 6), and the probability of buildings to become archaeological findings had decreased by 0.6 to 1 percent (columns 4 and 8). 

% Notes flat pre‑trend then sharp decline; explains SEs include Monte‑Carlo uncertainty.
For all these estimates it is seen, that there is a rather stable effect size before 1086-1208 of approximately zero, but after this the probability sharply declines. Appendix \ref{all_arch_param} provides tables of all parameters. It is worth considering whether these results can be spurious because of the added layer of uncertainty added by the Monte Carlo method. However, the standard errors used take this added uncertainty into account. These standard errors are constructed by resampling the Monte Carlo samples in a clustered fashion. Figure \ref{fig:arch_reg_boot} shows the parameter distribution in 1350 given by the bootstrap samples.

\begin{figure}
    \centering
    \caption{Archaelogical results}
    \begin{subfigure}[b]{0.45\textwidth}
        \centering
        \caption{Coins: Market access approach} \label{fig:arch1a}
        \includegraphics[width=\textwidth]{Plots/Regression_plots/arch_MA_coins.png}
    \end{subfigure}
    \hfill
    \begin{subfigure}[b]{0.45\textwidth}
        \centering
        \caption{Coins: Dummy approach} \label{fig:arch1b}
        \includegraphics[width=\textwidth]{Plots/Regression_plots/arch_dummy_coins.png}
    \end{subfigure}
    \vspace{0.45cm}
    \begin{subfigure}[b]{0.45\textwidth}
        \centering
        \caption{Buildings: Market access approach} \label{fig:arch1c}
        \includegraphics[width=\textwidth]{Plots/Regression_plots/arch_MA_buildings.png}
    \end{subfigure}
    \hfill
    \begin{subfigure}[b]{0.45\textwidth}
        \centering
        \caption{Buildings: Dummy approach} \label{fig:arch1d}
        \includegraphics[width=\textwidth]{Plots/Regression_plots/arch_dummy_buildings.png}
    \end{subfigure}
    \parbox{0.9\textwidth}{
    
    
    \caption*{\footnotesize \textit{Notes:} This shows 1000 draws from the bootstrap procedure, which takes classical (clustered) statistical uncertainty as well as dating uncertainty into account. Panel (a) shows the distribution of the effect on coin using the market access approach. Panel (b) shows the effect on coins using the dummy approach. Panel (c) shows the effect on buildings using the market access appraoch. Finally panel (d) shows the effect on buildings using the dummy approach. The dotted line is the point estimate.  \\ \textit{Source: Danish registry of archaeological findings}}
}    \label{fig:arch_reg}
\end{figure}


\begin{figure}
    \centering
    \caption{Distribution of parameter estimates in 1350}
    \begin{subfigure}[b]{0.45\textwidth}
        \centering
        \caption{Coins: Market access approach} \label{fig:distri_a}
        \includegraphics[width=\textwidth]{Plots/Regression_plots/arch_MA_coins_boot.png}
    \end{subfigure}
    \hfill
    \begin{subfigure}[b]{0.45\textwidth}
        \centering
        \caption{Coins: Dummy approach} \label{fig:distri_b}
        \includegraphics[width=\textwidth]{Plots/Regression_plots/arch_dummy_coins_boot.png}
    \end{subfigure}
    \vspace{0.45cm}
    \begin{subfigure}[b]{0.45\textwidth}
        \centering
        \caption{Buildings: Market access approach} \label{fig:distri_c}
        \includegraphics[width=\textwidth]{Plots/Regression_plots/arch_MA_buildings_boot.png}
    \end{subfigure}
    \hfill
    \begin{subfigure}[b]{0.45\textwidth}
        \centering
        \caption{Buildings: Dummy approach} \label{fig:distri_d}
        \includegraphics[width=\textwidth]{Plots/Regression_plots/arch_dummy_buildings_boot.png}
    \end{subfigure}
    \parbox{1\textwidth}{
\caption*{\footnotesize \textit{Notes:} This shows the archaeological regression results. Clustered bootstrapped standard errors are shown in the parentheses. Clustered at the parish level and sampled from the Monte Carlo procedure samples. Columns 1-4 show results using the full sample of all of Denmark. Columns 5-8 show results for a matched sample. All the even columns show results using the dummy definition of being affected. All the uneven columns show results using the change in market access approach. The outcome is the probability that a given finding type was generated in the area covered by that parish within $\pm$25 years of the reported year. A parameter estimate for all years 750, 800, ..., 1500 can be found in the Appendix \ref{all_arch_param}. \\ \textit{Source: Danish registry of archaeological findings}}
}
    \label{fig:arch_reg_boot}
\end{figure}

\begin{landscape}
\begin{table}
\centering
\caption{Archaeological regression results} \label{tab:arch1}
\footnotesize
\begin{tabular}{lcccccccc}
   \tabularnewline \midrule \midrule
                                                    & \multicolumn{4}{c}{Full sample}                                                    & \multicolumn{4}{c}{Matched sample} \\
   \midrule
   Outcome:                                         & \multicolumn{2}{c}{Coin findings} & \multicolumn{2}{c}{Buildings}                  & \multicolumn{2}{c}{Coin findings}               & \multicolumn{2}{c}{Buildings}\\
                                                    & (1)             & (2)             & (3)                   & (4)                    & (5)                    & (6)                    & (7)            & (8)\\  
                                                    & MA              & Dummy           & MA                    & Dummy                  & MA                     & Dummy                  & MA             & Dummy \\ 
   \midrule
   Year950 $\times$ Affected                        & -0.0042         & -0.0002         & 0.0011                & 0.0002                 & -0.0177$^{*}$          & -0.0021$^{*}$          & 0.0008          & $-8.62\times 10^{-6}$\\    
                                                    & (0.0038)        & (0.0004)        & (0.0018)              & (0.0003)               & (0.0096)               & (0.0011)               & (0.0039)        & (0.0006)\\   
   Year1050 $\times$ Affected                       & 0.0059$^{**}$   & 0.0008$^{***}$  & -0.0063$^{**}$        & -0.0007                & 0.0051                 & 0.0006                 & -0.0083         & -0.0012\\   
                                                    & (0.0025)        & (0.0003)        & (0.0032)              & (0.0006)               & (0.0039)               & (0.0005)               & (0.0064)        & (0.0010)\\   
   Year1150 $\times$ Affected                       & 0.0075          & 0.0003          & -0.0681$^{***}$       & -0.0065$^{*}$          & 0.0320                 & 0.0036                 & -0.0705$^{**}$  & -0.0078\\   
                                                    & (0.0121)        & (0.0011)        & (0.0170)              & (0.0033)               & (0.0245)               & (0.0029)               & (0.0318)        & (0.0051)\\   
   Year1250 $\times$ Affected                       & -0.0941$^{***}$ & -0.0120$^{***}$ & -0.0772$^{***}$       & -0.0066$^{*}$          & -0.1216$^{***}$        & -0.0150$^{***}$        & -0.0829$^{***}$ & -0.0093$^{*}$\\   
                                                    & (0.0199)        & (0.0018)        & (0.0193)              & (0.0038)               & (0.0350)               & (0.0045)               & (0.0318)        & (0.0054)\\   
   Year1350 $\times$ Affected                       & -0.1923$^{***}$ & -0.0193$^{***}$ & -0.0612$^{***}$       & -0.0065$^{**}$         & -0.1408$^{***}$        & -0.0159$^{***}$        & -0.0669$^{**}$  & -0.0076\\   
                                                    & (0.0273)        & (0.0037)        & (0.0185)              & (0.0031)               & (0.0414)               & (0.0058)               & (0.0297)        & (0.0046)\\   
   \midrule
   Observations                                     & 29,568          & 29,568          & 29,568                & 29,568                 & 7,424                  & 7,424                  & 7,424          & 7,424\\  
   Parishes                                         & 1848            & 1848            & 1848                  & 1848                   & 464                    & 464                    & 464            & 464\\
   \midrule 
   Parish FE                                        & Yes             & Yes             & Yes                   & Yes                    & Yes                    & Yes                    & Yes            & Yes\\  
   Year FE                                          & Yes             & Yes             & Yes                   & Yes                    & Yes                    & Yes                    & Yes            & Yes\\
   \midrule \midrule
\end{tabular}
\parbox{1.3\textwidth}{
\caption*{\footnotesize \textit{Notes:} Results from regressions on the probability that a given finding type was generated within $\pm$25 years in a parish. Standard errors (in parentheses) are clustered at the parish level and derived via bootstrapping of Monte Carlo samples. Columns 1-4 use the full sample and columns 5-8 use a matched sample; even columns use the dummy measure and odd columns the change in market access approach. Estimates for all years (750, 800, …, 1500) are provided in Appendix \ref{all_arch_param}. *** $p<0.01$, ** $p<0.05$, * $p<0.10$. \\ \textit{Source: Danish registry of archaeological findings}}
}
\end{table}
\end{landscape}


% Concludes negative effects corroborate main finding; similar magnitude vs. 1834 channel.
Whichever methodological alternatives are employed, the results show a large negative effect of the closure of the channel for the affected parishes. The results suggest that the closure of the channel between 1086 and 1208 had a significant impact on the West Limfjord regions, causing a decline in the probability of coin findings and buildings being generated in those areas. These findings are consistent with the hypothesis that the channel closure caused economic and commercial decline in the region. It is not possible to do an apples-to-apples comparison to the census data for the Agger channel of 1825. But it is noteworthy that the parameters are of similar sign and magnitudes for these two channels in the same location. This corroborates the external validity of the main result: First-nature geography is the cause of the location of prosperity across two very different societies. 

\FloatBarrier
\section{Conclusion}
% Summarizes main empirical findings: Agger channel → +26.7% population via market access, fertility 
This paper demonstrates that first–nature geography, particularly waterways, plays a critical role in determining the location of economic activity by shaping market access and that these forces act beyond being seeds of initial development. The unexpected emergence of the Agger channel in Denmark’s West Limfjord region in 1825 not only transformed trade but also catalyzed lasting economic change. The empirical analysis, based on a difference–in–differences design, shows that the channel’s creation led to a 26.7 percent increase in the population of affected parishes by 1901 — a growth driven primarily by enhanced market access that spurred trade, fishing, and manufacturing, and was fueled by higher fertility rather than migration.

% Highlights medieval reverse experiment; similar negative effects validate universality of geomorphology.
Furthermore, the reverse natural experiment from the medieval period — when a similar channel closed — produced analogous declines in economic activity, as evidenced by archaeological evidence. This external validation across two distinctly different societies, separated by 700 years, underscores the universal importance of geomorphology in shaping prosperity.

% Connects findings to modern challenges (climate change, Northwest Passage, flooding) and future research avenues.
In today’s rapidly changing world, where climate change, technological advances, and even space exploration are redefining physical and economic landscapes, understanding how first–nature characteristics influence economic outcomes is more relevant than ever. Just as the Agger channel reshaped an entire region, modern phenomena like the opening of the Northwest Passage near Greenland due to climate change, or the challenges of urban flooding, remind us that economic destiny is inextricably linked to the natural environment, leaving humanity as much as "prisoners of geomorphology" as it is agents of change. The approach introduced in this paper opens the door for further studies — such as those on the medieval channel in Bruges, Belgium \citep{Houtte1966, Charlier2011} or the historical dynamics of Königsberg/Kaliningrad \citep{Britannica2018} — to explore how changes in natural geography continue to influence economic trajectories.

% Final takeaway: shifting natural features shape prosperity; relevance for econ‑development debates.
In sum, this work contributes to a deeper understanding of how shifting natural features can unlock or constrain prosperity, offering insights that are both historically grounded and highly pertinent to contemporary debates on economic development and environmental change.

\newpage
\bibliographystyle{apacite}
\bibliography{main}

\normalsize
\newpage
\include{appendix}

\end{document}



